\subsection{Conclusion}
  \label{subsec:conclusion}
  In conclusion, we have shown that violations of Bell inequalities do not require
  dynamical nonlocality, nor the abandonment of realism, but can instead be understood
  as a generic consequence of non-injective projection.
  Within this framework, Bell-type correlations arise not from superluminal influences
  or hidden causal mechanisms, but from structural features of effective descriptions
  in which multiple underlying configurations correspond to the same observable outcomes.

  This perspective suggests that spacetime and its associated observables should not
  be regarded as the fundamental arena of physical reality, but rather as an effective
  level of description emerging from a richer underlying structure.
  Quantum correlations then cease to appear as instances of mysterious ``action at a
  distance'' and instead reflect global consistency constraints imposed by the
  non-injective mapping between underlying configurations—possibly of infra-physical
  or pre-geometric nature—and the domain of localized measurements.

  From this viewpoint, the empirical success of relativistic locality and quantum
  statistics is preserved at the effective level, while their apparent tension is resolved
  by recognizing the descriptive limitations inherent in any projected representation
  of a deeper relational substrate.
